% !PRINTED 604
\weberpage{604}

Diese Frage hat eine ganz allgemeine Beantwortung gefunden in einer Abhandlung von \gesperrt{Hilbert}\footnote{Ueber die Irreducibilit\"at ganzer rationaler Functionen mit ganzzahligen Co\"efficienten. Journal f\"ur Mathematik, Bd. 110.}, wo der allgemeine Satz bewiesen ist, dass man in einer irreduciblen Function beliebig vieler Variablen f\"ur einen beliebigen Theil der Variablen solche rationale Zahlen setzen kann, dass eine irreducible Function der \"ubrigen Variablen entsteht. Auf diesen allgemeinen Satz k\"onnen wir hier nicht eingehen.

Wir werden aber die gestellte Frage viel einfacher, wenn auch bei Weitem nicht so allgemein beantworten, indem wir zeigen, dass sich f\"ur jeden Primzahlgrad Gleichungen ohne Affect finden lassen.

Wir haben in \S.\,153, 9. bewiesen, dass eine transitive Permutationsgruppe von $n$ Ziffern, die nicht die symmetrische Gruppe ist, wenn $n$ eine Primzahl ist, keine einzelne Transposition enthalten kann.

Unter einer Gleichung mit einem Affect haben wir eine solche verstanden, deren Gruppe nicht die symmetrische ist. Hat also die irreducible Gleichung $f(x) = 0$ einen Affect und ist ihr Grad eine Primzahl, so muss ihre Gruppe $P$ transitiv sein, und sie kann keine Transposition zweier Wurzeln enthalten. Wenn wir also irgend $n - 2$ der Wurzeln dem Rationalit\"atsbereich adjungiren, so muss sie sich auf die Einheitsgruppe reduciren, da ja ausserdem nur noch die Vertauschung der beiden letzten Wurzeln \"ubrig bleiben k\"onnte, die in $P$ nicht vorkommt. Daraus folgt also, dass die beiden letzten Wurzeln in dem erweiterten K\"orper $\Omega$ enthalten sind, oder der Satz:

1. \gesperrt{Wenn eine irreducible Gleichung, deren Grad $n$ eine Primzahl ist, einen Affect hat, so k\"onnen zwei beliebige von ihren Wurzeln rational durch die \"ubrigen ausgedr\"uckt werden.}

Daraus folgt als Corollar:

2. \gesperrt{Wenn der K\"orper $\Omega$ nur reelle Zahlen enth\"alt, so kann eine irreducible Gleichung von Primzahlgrad $n$ mit einem Affect in $\Omega$ nicht zwei imagin\"are und $n - 2$ reelle Wurzeln haben.}

Nun giebt es aber unz\"ahlige Gleichungen von jedem be-
%CONT>>
